\section{Introduction}
The emergence of deep generative models~\cite{rombach2022high,ho2020denoising,takagi2023high,chen2023seeing,scotti2023reconstructing,lu2023minddiffuser,sun2024contrast,shen2024neuro,xia2024dream} has effectively bridged the modality gap between functional Magnetic Resonance Imaging (fMRI)~\cite{ogawa1990oxygenation} and diverse signals like visual, language, and audio.
However, many studies~\cite{,bai2023dreamdiffusion,lan2023seeing,zeng2023controllable,vanrullen2019reconstructing,dado2022hyperrealistic,shen2019deep,seeliger2018generative,NEURIPS2022_bee5125b,gu2022decoding,liu2024mind} focus on recovering the stimulated signal from brain activity, while interaction with reality requires active communication with human brain signals.

In this work, we explore a novel setting as shown in Fig.~\ref{fig:teaser} -- \emph{Can ``dreams'' be connected and actively influenced?} Such capability will present a novel and convenient interaction paradigm among people, facilitating numerous applications such as creative design via brainstorming.
The essence of this capability lies in steering the human ``dreams'' towards desired directions, ultimately \emph{enabling concept manipulation through direct operation on the fMRI signal}.
This paper mainly focuses on human ``dreams'' as a visual modality, but the fMRI signal can also store other types of stimulated information, such as audio~\cite{defossez2023decoding} or language~\cite{zhao2024mapguide}. Consequently, our proposed approach can be seamlessly extended to these scenarios.


This task requires direct instruction of fMRI signals, which presents notable challenges on two fronts: 1) The inherent abstractness and ambiguity of collected brain signals often impede the comprehension of their content;
2) A significant modality gap exists between fMRI signals and natural language instructions, hindering the learning of accurate associations between brain activity features and language instructions. Since not all the information within the fMRI signals is relevant to the intended instruction, successful instruction requires \emph{pinpointing and modifying the relevant features while preserving the irrelevant ones}.



We introduce \textbf{\model{}}, a dual-stream diffusion framework tailored for fMRI signal instruction to address these challenges. The key is to \emph{effectively guide language instructions to hone in on and modify the pertinent information embedded within fMRI signals}. To establish a correlation between these distinct signals, we first employ a Stable Diffusion (SD)~\cite{rombach2022high} backbone to interpret brain activity using visual prior knowledge. Then, we integrate language instructions into a parallel SD backbone to manipulate the acquired visual priors.
This parallel forwarding approach has proven to be effective in improving network learning ability by leveraging the homogeneous nature of the extracted features~\cite{zhang2023adding}. To bridge these two streams, i.e., the interpretation stream and instruction stream, we introduce an adaptor network designed to modulate the intermediate visual contents from the interpretation stream hierarchically.

Ideally, with this design, the framework will operate under a \emph{progressive interpretation and instruction paradigm, concurrently identifying and manipulating correlated regions}. However, the diffusion operation behaves distinctively along the temporal dimension, synthesizing varied content across different time steps~\cite{zhang2023prospect,Zhang_2023_inst}. At specific time steps, the second stream must identify specific visual features for instruction, necessitating that the first stream completes their formation beforehand. Therefore, we devise an asynchronous diffusion strategy to ensure that the first stream has adequate time to develop overall semantics, thereby facilitating smooth manipulation progress. Additionally, to further guide the model toward intended spatial locations, we introduce a Large Language Model (LLM)~\cite{wei2022emergent} agent tasked with identifying regions relevant to instructions. Using the identified areas, we implement a region-aware attention strategy to ensure that the instructions are applied precisely within those regions.
\textbf{Our main contributions can be summarized as:}

\noindent \textbf{(1)} We propose a visual brainstorming instruction system, named \textbf{\model{}}, empowering users to effectively interfacing with human ``dreams'' to perform desired operation.

\noindent \textbf{(2)} We devise a dual-stream diffusion network coupled with an adaptor to seamlessly translate fMRI signals towards their intended directions.

 \noindent \textbf{(3)} We introduce an asynchronous diffusion strategy, enhanced with LLM-guided region-aware manipulation, to facilitate faithful instruction within pertinent regions from both temporal and spatial perspectives.
\noindent \textbf{(4)} While designed for language-guided instruction, our framework can readily adapt to visual instruction or multi-modal instruction with minimal adjustments.



